% Part: sets-functions-relations
% Chapter: relations
% Section: special-properties

\documentclass[../../../include/open-logic-section]{subfiles}

\begin{document}

\olfileid{sfr}{rel}{prp}
\olsection{સંબંધોના વિશેષ ગુણધર્મો}

\begin{intro}
કેટલાક પ્રકારના સંબંધો એટલા સામાન્ય છે કે તેમને વિશેષ નામ
આપવામાં આવ્યાં છે. દાખલા તરીકે, $\le$ અને $\subseteq$ બંને
પોતપોતાના પ્રદેશો પર (માનો કે $\le$ માટે $\Nat$ અને
$\subseteq$ માટે $\Pow{A}$) એકસરખી રીતે સંબંધ સ્થાપે છે.
આ સંબંધો કઈ ચોક્કસ રીતે સરખા છે અને કઈ રીતે જુદા પડે છે તે
સમજવા માટે, સંબંધોમાં હોઈ શકે એવા કેટલાક વિશેષ ગુણધર્મો મુજબ
આપણે તેમનું વર્ગીકરણ કરીએ છીએ. આ વિશેષ ગુણધર્મોમાંના કેટલાક
(અથવા તેમના સંયોજનો) ખાસ મહત્વના નીવડે છે: ક્રમો અને સામ્ય સંબંધો.
\end{intro}

\begin{defn}[સ્વવાચકતા]
સંબંધ $R \subseteq A^2$ \emph{સ્વવાચક} હોય તો અને તો જ,
દરેક $x \in A$ માટે, $Rxx$.
\end{defn}

\begin{defn}[પરંપરિતતા]
સંબંધ $R \subseteq A^2$ \emph{પરંપરિત} હોય તો અને તો જ,
જ્યારે પણ $Rxy$ અને $Ryz$, ત્યારે $Rxz$ પણ.
\end{defn}

\begin{defn}[સંમિતતા]
સંબંધ $R \subseteq A^2$ \emph{સંમિત} હોય તો અને તો જ,
જ્યારે પણ $Rxy$, ત્યારે $Ryx$ પણ.
\end{defn}

\begin{defn}[વિસંમિતતા]
સંબંધ $R \subseteq A^2$ \emph{વિસંમિત} હોય તો અને તો જ,
જ્યારે પણ $Rxy$ અને $Ryx$ બંને, ત્યારે $x=y$
(અથવા, બીજા શબ્દોમાં: જો $x\neq y$, તો
કાં $\lnot Rxy$ અથવા $\lnot Ryx$).
\end{defn}

\begin{explain}
સંમિત સંબંધમાં $Rxy$ અને $Ryx$ હંમેશાં એકસાથે સાચાં હોય છે,
અથવા બંનેમાંથી એકેય સાચું હોતું નથી. વિસંમિત સંબંધમાં
$Rxy$ અને $Ryx$ એકસાથે સાચાં હોય, તો તે માત્ર $x = y$
હોય ત્યારે જ શક્ય છે. નોંધો કે આમાં $x = y$ હોય ત્યારે
$Rxy$ અને $Ryx$ સાચાં હોવાં \emph{જોઈએ} એવી માંગ નથી;
માત્ર એ શક્યતા બાકાત રખાતી નથી. તેથી વિસંમિત સંબંધ
સ્વવાચક હોઈ શકે છે, પણ દરેક વિસંમિત સંબંધ સ્વવાચક હોય
એવું નથી. એ પણ નોંધો કે વિસંમિત હોવું અને માત્ર સંમિત
ન હોવું એ જુદી શરતો છે. ખરેખર, કોઈ સંબંધ એકસાથે સંમિત
અને વિસંમિત બંને હોઈ શકે છે (દા.ત., તાદાત્મ્ય સંબંધ એવો છે).
\end{explain}

\begin{defn}[તુલનાયુક્તતા]
સંબંધ $R \subseteq A^2$ \emph{તુલનાયુક્ત} કહેવાય જો,
દરેક $x,y\in A$ માટે, $x \neq y$ હોય તો
કાં $Rxy$ અથવા $Ryx$.
\end{defn}

\begin{prob}
એવા સંબંધોના દાખલા આપો જે (a) સ્વવાચક અને સંમિત હોય પણ
પરંપરિત ન હોય, (b) સ્વવાચક અને વિસંમિત હોય,
(c) વિસંમિત અને પરંપરિત હોય પણ સ્વવાચક ન હોય, અને
(d) સ્વવાચક, સંમિત અને પરંપરિત હોય.
સંખ્યાઓ કે ગણો પરના સંબંધો વાપરશો નહીં.
\end{prob}

\begin{defn}[અસ્વવાચકતા]
સંબંધ $R \subseteq A^2$ \emph{અસ્વવાચક} કહેવાય જો,
દરેક $x \in A$ માટે, $Rxx$ સાચું ન હોય.
\end{defn}

\begin{defn}[અસંમિતતા]
સંબંધ $R \subseteq A^2$ \emph{અસંમિત} કહેવાય જો,
કોઈ પણ જોડ $x,y\in A$ માટે $Rxy$ અને $Ryx$ બંને સાચાં ન હોય.
\end{defn}

નોંધો કે જો $A \neq \emptyset$, તો $A$ પરનો કોઈ અસ્વવાચક
સંબંધ સ્વવાચક નથી અને $A$ પરનો દરેક અસંમિત સંબંધ વિસંમિત
પણ છે. જોકે, એવા $R \subseteq A^2$ છે જે સ્વવાચક પણ
નથી અને અસ્વવાચક પણ નથી, અને એવા વિસંમિત સંબંધો છે
જે અસંમિત નથી.

\end{document}
